\chapter{2009年山东卷（文）}

\section{单选题}

\begin{problem}
集合 \(A = \{0, 2, a\}\)，\(B = \{1, a^2\}\). 若 \(A \cup B = \{0, 1, 2, 4, 16\}\)，则 \(a\) 的值为（\quad）

\choices
    {0}
    {1}
    {2}
    {4}
\end{problem}

\begin{problem}
复数 \(\frac{3 - \i}{1 - \i}\) 等于（\quad）

\choices
    {\(1 + 2\i\)}
    {\(1 - 2\i\)}
    {\(2 + \i\)}
    {\(2 - \i\)}
\end{problem}

\begin{problem}
将函数 \( y = \sin 2x \) 的图象向左平移 \( \frac{\pi}{4} \) 个单位，再向上平移1个单位，所得图象的函数解析式是（\quad）

\choices
    {\(y = \cos 2x\)}
    {\(y = 2\cos^2 x\)}
    {\(y = 1 + \sin \left(2x + \frac{\pi}{4}\right)\)}
    {\(y = 2\sin^2 x\)}
\end{problem}

\begin{problem}
一空间几何体的三视图如图所示，则该几何体的体积为

\bitmapfigure[max width=0.56\linewidth,max totalheight=6cm]{img/2009/shandong_liberal/q4_fig1.png}

\choices
    {\(2\pi + 2\sqrt{3}\)}
    {\(4\pi + 2\sqrt{3}\)}
    {\(2\pi + \frac{2\sqrt{3}}{3}\)}
    {\(4\pi + \frac{2\sqrt{3}}{3}\)}
\end{problem}

\begin{problem}
在 \(\R\) 上定义运算 \(\odot: a \odot b = ab + 2a + b\)，则满足 \(x \odot (x - 2) < 0\) 的实数 \(x\) 的取值范围为（\quad）

\choices
    {(0,2)}
    {\((-2,1)\)}
    {\((- \infty, -2) \cup (1, +\infty)\)}
    {\((-1, 2)\)}
\end{problem}

\begin{problem}
函数 \( y = \frac{\e^x + \e^{-x}}{\e^x - \e^{-x}} \) 的图象大致为（\quad）

\bitmapfigure[max width=0.88\linewidth,max totalheight=6.2cm]{img/2009/shandong_liberal/q6_options.png}
\end{problem}

\begin{problem}
定义在 \(\R\) 上的函数 \(f(x)\) 满足 \(f(x) = \begin{cases} \log_2(4 - x), & x \leqslant 0, \\ f(x - 1) - f(x - 2), & x > 0. \end{cases}\) 则 \(f(3)\) 的值为（\quad）

\choices
    {\(-1\)}
    {\(-2\)}
    {\(1\)}
    {\(2\)}
\end{problem}

\begin{problem}
设 \(P\) 是 \(\triangle ABC\) 所在平面内的一点，\(\overrightarrow{BC} + \overrightarrow{BA} = 2\overrightarrow{BP}\)，则（\quad）

\choices
    {\(\overrightarrow{PA} +\overrightarrow{PB} = 0\)}
    {\(\overrightarrow{PB} +\overrightarrow{PC} = 0\)}
    {\(\overrightarrow{PC} +\overrightarrow{PA} = 0\)}
    {\(\overrightarrow{PA} +\overrightarrow{PB} +\overrightarrow{PC} = 0\)}
\end{problem}

\begin{problem}
已知 \(\alpha, \beta\) 表示两个不同的平面，\(m\) 为平面 \(\alpha\) 内的一条直线，则“\(\alpha \perp \beta\)”是“\(m \perp \beta\)”的（\quad）

\choices
    {充分不必要条件}
    {必要不充分条件}
    {充要条件}
    {既不充分也不必要条件}
\end{problem}

\begin{problem}
设斜率为 2 的直线 l 过抛物线 \( y^{2} = ax (a \neq 0) \) 的焦点 F，且和 y 轴交于点 A. 若 \( \triangle OAF \) (O 为坐标原点) 的面积为 4，则抛物线方程为（\quad）

\choices
    {\( y^{2} = \pm 4x \)}
    {\( y^{2} = \pm 8x \)}
    {\( y^{2} = 4x \)}
    {\( y^{2} = 8x \)}
\end{problem}

\begin{problem}
在区间 \(\left[-\frac{\pi}{2}, \frac{\pi}{2}\right]\) 上随机取一个数 \(x\)，\(\cos x\) 的值介于 0 到 \(\frac{1}{2}\) 之间的概率为（\quad）

\choices
    {\( \frac{1}{3} \)}
    {\( \frac{2}{\pi} \)}
    {\( \frac{1}{2} \)}
    {\( \frac{2}{3} \)}
\end{problem}

\begin{problem}
已知定义在 \(\R\) 上的奇函数 \(f(x)\) 满足 \(f(x - 4) = -f(x)\). 且在区间 \([0,2]\) 上是增函数，则（\quad）

\choices
    {\(f(-25) < f(11) < f(80)\)}
    {\(f(80) < f(11) < f(-25)\)}
    {\(f(11) < f(80) < f(-25)\)}
    {\(f(-25) < f(80) < f(11)\)}
\end{problem}

\section{填空题}

\begin{problem}
在等差数列 \(\{a_{n}\}\) 中，\(a_{3} = 7, a_{5} = a_{2} + 6\)，则 \(a_{6} = \)\fillinblank{}.
\end{problem}

\begin{problem}
若函数 \(f(x) = a^x - x - a\)（\(a > 0\) 且 \(a \neq 1\)）有两个零点，则实数 \(a\) 的取值范围是\fillinblank{}.
\end{problem}

\begin{problem}
执行如图的程序框图，输出的 \(T =\)\fillinblank{}.

\bitmapfigure[max width=0.42\linewidth,max totalheight=5.8cm]{img/2009/shandong_liberal/q15_fig1.png}
\end{problem}

\begin{problem}
某公司租赁甲，乙两种设备生产 A, B 两类产品，甲种设备每天能生产 A 类产品 5 件和 B 类产品 10 件，乙种设备每天能生产 A 类产品 6 件和 B 类产品 20 件. 已知设备甲每天的租赁费为 200 元，设备乙每天的租赁费为 300 元. 现该公司至少要生产 A 类产品 50 件，B 类产品 140 件，所需租赁费最少为\fillinblank{}\text{元}.
\end{problem}

\section{解答题}

\begin{problem}
已知函数 \(f(x) = 2\sin x\cos^2\frac{x}{2} + \cos x\sin \varphi - \sin x\)（\(0 < \varphi < \pi\)）在 \(x = \pi\) 处取最小值.

\begin{enumerate}
\item 求 \(\varphi\) 的值.
\item 在 \(\triangle ABC\) 中，\(a, b, c\) 分别是角 \(A, B, C\) 的对边. 已知 \(a = 1\)，\(b = \sqrt{2}\)，\(f(A) = \frac{\sqrt{3}}{3}\)，求角 \(C\).
\end{enumerate}
\end{problem}

\begin{problem}
如图，在直四棱柱 \(ABCD - A_{1}B_{1}C_{1}D_{1}\) 中，底面 \(ABCD\) 为等腰梯形，\(AB // CD, AB = 4, BC = CD = 2, AA_{1} = 2, E, E_{1}\) 分别是棱 \(AD, AA_{1}\) 的中点.

\begin{enumerate}
\item 设 F 是棱 AB 的中点，证明：直线 \(EE_{1} //\) 平面 \(FCC_{1}\)；
\item 证明：平面 \(D_{1}AC \perp\) 平面 \(BB_{1}C_{1}C\).
\end{enumerate}

\bitmapfigure[max width=0.52\linewidth,max totalheight=4.8cm]{img/2009/shandong_liberal/q18_fig1.png}

\end{problem}

\begin{problem}
一汽车厂生产 \(A, B, C\) 三类轿车，每类轿车均有舒适型和标准型两种型号，某月的产量如下表（单位：辆）：

\begin{center}
\begin{tabular}{|c|c|c|c|}
\hline
& 轿车 A & 轿车 B & 轿车 C \\
\hline
舒适型 & 100 & 150 & \(z\) \\
\hline
标准型 & 300 & 450 & 600 \\
\hline
\end{tabular}
\end{center}

按类型分层抽样的方法在这个月生产的轿车中抽取 50 辆，其中有 A 类轿车 10 辆.

\begin{enumerate}
\item 求 \(z\) 的值.
\item 用分层抽样的方法在 C 类轿车中抽取一个容量为 5 的样本. 将该样本看成一个总体，从中任取 2 辆，求至少有 1 辆舒适型轿车的概率.
\item 用随机抽样的方法从 B 类舒适型轿车中抽取 8 辆，经检测它们的得分如下：9.4, 8.6, 9.2, 9.6, 8.7, 9.3, 9.0, 8.2. 把这 8 辆轿车的得分看作一个总体，从中任取一个数，求该数与样本平均数之差的绝对值不超过 0.5 的概率.
\end{enumerate}
\end{problem}

\begin{problem}
等比数列 \(\{a_{n}\}\) 的前 \(n\) 项和为 \(S_{n}\). 已知对任意的 \(n \in \mathbf{N}^{*}\)，点 \((n, S_{n})\) 均在函数 \(y = b^{x} + r\)（\(b > 0\) 且 \(b \neq 1\)，\(b, r\) 均为常数）的图象上.

\begin{enumerate}
\item 求 \(r\) 的值.
\item 当 \(b = 2\) 时，记 \(b_{n} = \frac{n + 1}{4a_{n}}\)（\(n \in \mathbf{N}^{*}\)），求数列 \(\{b_{n}\}\) 的前 \(n\) 项和 \(T_{n}\).
\end{enumerate}
\end{problem}

\begin{problem}
设 \(m \in \R\)，在平面直角坐标系中，已知向量 \(a = (mx, y + 1)\)，向量 \(b = (x, y - 1)\)，\(a \perp b\)，动点 \(M(x, y)\) 的轨迹为 \(E\).

\begin{enumerate}
\item 求轨迹 \(E\) 的方程，并说明该方程所表示曲线的形状.
\item 已知 \(m = \frac{1}{4}\)，证明：存在圆心在原点的圆，使得该圆的任意一条切线与轨迹 \(E\) 恒有两个交点 \(A, B\)，且 \(OA \perp OB\)（\(O\) 为坐标原点），并求出该圆的方程.
\item 已知 \(m = \frac{1}{4}\)，设直线 \(l\) 与圆 \(C: x^2 + y^2 = R^2\)（\(1 < R < 2\)）相切于 \(A_1\)，且 \(l\) 与轨迹 \(E\) 只有一个公共点 \(B_1\)，当 \(R\) 为何值时，\(|A_1B_1|\) 取得最大值? 并求最大值.
\end{enumerate}
\end{problem}

\begin{problem}
已知函数 \(f(x) = \frac{1}{3}ax^3 + bx^2 + x + 3\)，其中 \(a \neq 0\).

\begin{enumerate}
\item 当 \(a, b\) 满足什么条件时，\(f(x)\) 取得极值?
\item 已知 \(a > 0\)，且 \(f(x)\) 在区间 \((0,1]\) 上单调递增，试用 \(a\) 表示出 \(b\) 的取值范围.
\end{enumerate}
\end{problem}
